\section*{Chapter \ref{ch:dv:local-volatility} --- Local Volatility}

\begin{solution}{exo:dv:local-volatility:1}
$\sigma_{\mathrm{loc}}=\sqrt{0.045/0.9}=\sqrt{0.05}=22.4\%$.
\end{solution}

\begin{solution}{exo:dv:local-volatility:2}
$16.4\%+2\times(-0.47)\times(-0.1)=25.7\%$, against 27.4\%: the rule captures the slope; the rest is the
smile's curvature, which the \omterm{def:dv:local-volatility:model}{local volatility} also doubles.
\end{solution}

\begin{solution}{exo:dv:local-volatility:3}
Independent of the spot: $\sigma_{\mathrm{loc}}^2=0.5\times0.15^2+0.5\times0.30^2=0.05625$, a flat 23.7\%. At 90
with probability 0.8 of the high state: $0.8\times0.09+0.2\times0.0225=0.0765$, 27.7\%.
\end{solution}

\begin{solution}{exo:dv:local-volatility:4}
The formula divides the calendar spread by the butterfly: a \omterm{def:dv:implied-volatility-and-its-surface:static}{calendar arbitrage} makes the local variance
negative, a \omterm{def:dv:implied-volatility-and-its-surface:static}{butterfly arbitrage} makes the denominator negative or zero, and either gives no \omterm{def:dv:local-volatility:model}{local volatility}.
A noisy quote bends the smile locally, and the second derivative in the denominator amplifies that bend into a
spike or a hole in the \omterm{def:dv:local-volatility:model}{local volatility} near the quote.
\end{solution}

\begin{solution}{exo:dv:local-volatility:5}
$\sqrt{(0.1919^2\times1-0.1779^2\times0.5)/0.5}=20.49\%$, against 20.68\% for the model's forward-start at-the-money
option: the level of the \omterm{def:dv:local-volatility:forward}{forward smile} is essentially forward variance; the small gap is the curvature of the
\omterm{def:dv:local-volatility:forward}{forward smile} at the money.
\end{solution}

\begin{solution}{exo:dv:local-volatility:6}
$2\times0.47\times0.01=0.94$ point, against 1.31 simulated: the rule holds for short expiries and weak skews, and
this three-month skew is steep.
\end{solution}

\begin{solution}{exo:dv:local-volatility:7}
With 50\,000 paths the 90 put's error is $-0.04$ point with a standard error of 0.13; with 200\,000, $+0.11$ with a
standard error of 0.06. The errors of all strikes at 200\,000 paths have the same sign: a discretisation bias of the
daily time step (about 0.08 point), which more paths do not remove and four steps a day does.
\end{solution}

\begin{solution}{exo:dv:local-volatility:8}
A European payoff on the index at one date, yes: its price depends only on the marginal distribution at that date,
which the fit reproduces. A forward-start option pays on the ratio of the index at two dates, which depends on the
joint distribution, and so on the model's dynamics: \omterm{def:dv:local-volatility:model}{local volatility} gives it a flatter \omterm{def:dv:local-volatility:forward}{forward smile} than a
stochastic volatility model fitted to the same vanillas.
\end{solution}

\begin{solution}{pb:dv:local-volatility:1}
\textbf{1.} 16.4\% and 17.7\% at the money; 21.1\% and 27.4\% at $k=-0.1$.
\textbf{2.} 1.98.
\textbf{3.} Implied volatility averages the \omterm{def:dv:local-volatility:model}{local volatility} between the spot and the strike; to produce a given
slope of the average, the \omterm{def:dv:local-volatility:model}{local volatility} must slope twice as fast.
\textbf{4.} $+0.08$, $+0.06$, $+0.09$ volatility point.
\textbf{5.} A grid or extrapolation problem, or an arbitrage in the input surface, not noise.
\textbf{6.} 20.7\% and 17.8\%.
\textbf{7.} The forward at-the-money volatility from six months to one year is 20.5\% from the term structure: the
\omterm{def:dv:local-volatility:forward}{forward smile} sits at the forward variance.
\textbf{8.} $-0.19$ and $-0.34$ per unit of \omterm{def:dv:implied-volatility-and-its-surface:surface}{log-moneyness}.
\textbf{9.} 0.57.
\textbf{10.} Its \omterm{def:dv:local-volatility:model}{local volatility} is steep near today's spot at short times; six months out the spot has spread and
the relevant local volatilities are those of longer times, where the function is flatter.
\textbf{11.} 21.8\%; today's shape shifted to the forward level gives 22.5\%.
\textbf{12.} 3.82 and 4.00 per 100.
\textbf{13.} \omterm{def:dv:local-volatility:model}{Local volatility}, the cheaper model for the buyer's protection it sells; a model chosen because it prices
the risk lower is a model risk (chapter 27).
\textbf{14.} Forward-start options or cliquets quoted by dealers, or the variance of future implied skews observed
historically, which reveal how skews behave.
\textbf{15.} Price with both models and hold the difference, 0.18 per 100, as a model reserve until the risk is gone.
\textbf{16.} 23.4\% under \omterm{def:dv:local-volatility:model}{local volatility}, 18.8\% under \omterm{def:dv:implied-volatility-and-its-surface:sticky}{sticky strike}, from 16.4\%.
\textbf{17.} For products whose value depends on the marginal distributions and short-dated dynamics (barriers near
the spot, American options), and as the \omterm{def:dv:local-volatility:projection}{Markovian projection} that other models are calibrated to.
\textbf{18.} Randomness in volatility that is not a function of the spot: stochastic volatility (chapter 10) or a
forward-variance model (chapter 12).
\textbf{19.} 0.57: the six-month skew six months forward is 57\% of today's six-month skew.
\textbf{20.} The dynamics: how the smile moves with the spot and how it will look at future dates.
\end{solution}

\begin{solution}{iq:dv:local-volatility:1}
The volatility as a deterministic function of time and spot, $\sigma_{\mathrm{loc}}(t,S)$, of a diffusion that
reprices every European option. It is unique because the call prices determine the density at every date
(Breeden--Litzenberger), and the density's evolution determines the diffusion coefficient (Dupire's formula).
\iqlookfor{from prices to densities to the coefficient.}
\end{solution}

\begin{solution}{iq:dv:local-volatility:2}
It rises, by about twice what \omterm{def:dv:implied-volatility-and-its-surface:sticky}{sticky strike} gives: the at-the-money implied volatility is close to the \omterm{def:dv:local-volatility:model}{local
volatility} at the spot, which slopes twice as steeply as the implied skew, and the spot has moved to a region of
higher \omterm{def:dv:local-volatility:model}{local volatility}.
\iqlookfor{the factor two and its reason.}
\end{solution}

\begin{solution}{iq:dv:local-volatility:3}
Density $q$ solves $\partial_Tq=\tfrac12\partial_{ss}(\sigma^2s^2q)$; $C=\int(s-K)^+q\,ds$; $\partial_TC=
\tfrac12\int(s-K)^+\partial_{ss}(\sigma^2s^2q)\,ds=\tfrac12\sigma^2(T,K)K^2q(T,K)$ after two integrations by parts;
$q=\partial_{KK}C$. So $\sigma^2(T,K)=2\partial_TC/(K^2\partial_{KK}C)$.
\iqlookfor{the forward equation and the integrations by parts.}
\end{solution}

\begin{solution}{iq:dv:local-volatility:4}
They agree on anything that depends only on the distribution of the underlying at single dates (European payoffs,
including digitals and variance swaps without jumps). They can differ on anything that depends on the joint law at
several dates or on paths: barriers, Asian, lookback and forward-start options, cliquets, autocallables.
\iqlookfor{marginals against joint laws.}
\end{solution}

\begin{solution}{iq:dv:local-volatility:5}
From differentiating a noisy or arbitrageable surface, from interpolation kinks between slices, and from
extrapolation beyond the quoted strikes. Fit a smooth arbitrage-free parametric surface first, compute the \omterm{def:dv:local-volatility:model}{local
volatility} from it in total variance, cap and floor it where the input has no information, and check by repricing.
\iqlookfor{smooth input, not smoothing the output.}
\end{solution}

\begin{solution}{iq:dv:local-volatility:6}
The \omterm{def:dv:local-volatility:model}{local volatility model} with $\sigma^2(t,K)=\E[\sigma_t^2\mid S_t=K]$, which has the same marginals as the
stochastic volatility model. To fit a stochastic volatility model to the whole surface, multiply its volatility by a
leverage function $L(t,S)$ chosen so that $L^2(t,K)\,\E[\sigma_t^2\mid S_t=K]$ equals the market's local variance: the
stochastic-local volatility model of chapter 20.
\iqlookfor{conditional expectation of variance given spot, and the leverage function.}
\end{solution}
